% Full-page figure: model triangulation. Four independent routes to the
% same engineering quantity (UK personal-transport energy demand, after
% MacKay), with their values and intervals plotted on a common axis,
% showing both agreement and disagreement.
\begin{figure}[p]
\centering
\begin{tikzpicture}[font=\small]

% ===== Title block =====
\node[font=\large\bfseries, align=center] at (5.5, 13.0)
  {Four routes to one quantity: model triangulation};
\node[font=\small\itshape, align=center, color=engblack!70] at (5.5, 12.4)
  {Annual personal-transport energy demand per UK resident, current as of 2008-09.};
\node[font=\small\itshape, align=center, color=engblack!70] at (5.5, 12.0)
  {Four independent computation routes; the spread of their estimates is the structural-uncertainty floor.};

% ===== Route boxes =====

% Route 1
\node[draw=engblack, fill=engblue!10, rounded corners, minimum width=10cm,
  minimum height=1.9cm, align=left, anchor=north west, font=\footnotesize] at (0.5, 11.0)
  {\textbf{Route 1: vehicle-kilometre $\times$ energy intensity.} \\
   $E_1 = D \cdot \eta = (13{,}500\,\text{km}/\text{yr})(0.80\,\text{kWh}/\text{km})
   = 10{,}800\,\text{kWh}/\text{yr}.$\\
   $D$ from UK National Travel Survey 2009; $\eta$ from typical petrol-car
   energy-per-km (cars only). \\
   \textit{Uncertainty: $\pm 20\%$ from $\eta$ spread across vehicle fleet.}};

% Route 2
\node[draw=engblack, fill=engblue!10, rounded corners, minimum width=10cm,
  minimum height=1.9cm, align=left, anchor=north west, font=\footnotesize] at (0.5, 8.9)
  {\textbf{Route 2: fuel sales $\div$ population.}\\
   $E_2 = F / N = (44 \cdot 10^9\,\text{L petrol-equiv}/\text{yr})
   (9.7\,\text{kWh}/\text{L}) / (61 \cdot 10^6) = 7{,}000\,\text{kWh}/\text{yr}.$\\
   $F$ from UK Department for Transport fuel-sales tables; $N$ from ONS.\\
   \textit{Uncertainty: $\pm 5\%$ from sales aggregation; lower than Route 1.}};

% Route 3
\node[draw=engblack, fill=engblue!10, rounded corners, minimum width=10cm,
  minimum height=1.9cm, align=left, anchor=north west, font=\footnotesize] at (0.5, 6.8)
  {\textbf{Route 3: travel-time budget $\times$ average power.}\\
   $E_3 = t \cdot P = (1.1\,\text{h}/\text{day})(365)(15\,\text{kW}\,\text{avg per car})/4
     = 4{,}500\,\text{kWh}/\text{yr per person}.$\\
   $t$ from time-use surveys; $P$ from typical engine load; factor 4 occupancy correction.\\
   \textit{Uncertainty: $\pm 30\%$ from occupancy and power assumptions.}};

% Route 4
\node[draw=engblack, fill=engblue!10, rounded corners, minimum width=10cm,
  minimum height=1.9cm, align=left, anchor=north west, font=\footnotesize] at (0.5, 4.7)
  {\textbf{Route 4: CO\textsubscript{2} emissions $\div$ carbon intensity.}\\
   $E_4 = (M_\text{CO\textsubscript{2},transport} / I_\text{CO\textsubscript{2}}) / N
     \approx 8{,}000\,\text{kWh}/\text{yr}.$\\
   $M$ from UK national GHG inventory (transport sector minus freight);
   $I$ from fuel carbon content.\\
   \textit{Uncertainty: $\pm 15\%$ from sectoral attribution.}};

% ===== Common axis with intervals =====
\draw[->, thick, draw=engblack!70] (1, 2.5) -- (10, 2.5)
  node[right, font=\small] {$E$ (kWh/yr per resident)};
\foreach \xt/\xv in {2/2000, 3.5/4000, 5/6000, 6.5/8000, 8/10000, 9.5/12000} {
  \draw[draw=engblack!70] (\xt, 2.5) -- (\xt, 2.4);
  \node[font=\scriptsize, anchor=north] at (\xt, 2.35) {\xv};
}

% Interval bars on the axis
% Route 1: 10800 +/- 20% => bar from 8640 to 12960, x-coord = 1 + 9 * E/12000
% map: x = 1 + (E/12000)*9
\foreach \E/\sig/\col/\label in {
  10800/2160/engblue!85/{R1},
  7000/350/engblue!85/{R2},
  4500/1350/engblue!85/{R3},
  8000/1200/engblue!85/{R4}} {
  \pgfmathsetmacro\xc{1 + (\E/12000)*9}
  \pgfmathsetmacro\xl{1 + ((\E-\sig)/12000)*9}
  \pgfmathsetmacro\xr{1 + ((\E+\sig)/12000)*9}
  \draw[thick, draw=\col] (\xl, 1.7) -- (\xr, 1.7);
  \draw[thick, draw=\col] (\xl, 1.6) -- (\xl, 1.8);
  \draw[thick, draw=\col] (\xr, 1.6) -- (\xr, 1.8);
  \filldraw[\col] (\xc, 1.7) circle (0.08);
  \node[font=\scriptsize, anchor=south, color=\col] at (\xc, 1.85) {\label};
}

% Annotation: triangulation envelope
\draw[dashed, draw=engred!80, thick] (4.375, 1.3) -- (4.375, 2.0);
\draw[dashed, draw=engred!80, thick] (8.6, 1.3) -- (8.6, 2.0);
\node[font=\footnotesize, color=engred!80, align=center, anchor=north] at (6.5, 1.25)
  {triangulation envelope: $\sim$4{,}500 to $\sim$11{,}000 kWh/yr,\\
   the structural-uncertainty range the four-route comparison reports};

% Conclusion box
\node[draw=engblack!60, fill=engblack!5, rounded corners,
  minimum width=10cm, minimum height=1.1cm, align=left, anchor=north west,
  font=\footnotesize] at (0.5, 0.9)
  {\textbf{Reading.} The four routes agree on the order of magnitude
   ($\sim$$10^4$ kWh/yr per resident) and disagree on the factor-of-two
   level. The disagreement is the structural floor on any single-route
   estimate. A reader who reports only Route 2's tight $\pm 5\%$
   interval has under-reported the modelling uncertainty by a factor of
   five, because Route 2's interval is the within-model interval; the
   across-route interval is what the planner using the number needs.};

\end{tikzpicture}
\caption[Model triangulation: four routes to UK personal-transport
energy demand]{Four independent computation routes to a single
engineering quantity: annual personal-transport energy demand per
UK resident, current as of 2008-09. Each route's value and within-route
uncertainty are plotted on the common axis; the dashed red envelope
marks the spread across routes, which sets the lower bound on the
structural uncertainty of any single-route estimate. The across-route
spread is six times Route~2's reported within-route interval, which is
the chapter's point: a model whose interval is reported without
triangulation under-reports the uncertainty the decision-maker
actually faces. The example is drawn from David MacKay's
\emph{Sustainable Energy: Without the Hot Air}
\cite{gen:mackay2009}, which builds the triangulation discipline at
book length on energy questions.}
\label{fig:vol02:ch18:triangulation}
\end{figure}
